\documentclass{article}
\usepackage{lmodern}
\usepackage{amsmath}
\usepackage{listings}
\usepackage{fullpage}
\usepackage{parskip}
\usepackage{xcolor}
\lstset{
	basicstyle=\ttfamily,
	columns=fullflexible,
	frame=single,
	breaklines=true,
	postbreak=\mbox{\textcolor{red}{$\hookrightarrow$}\space},
}
\begin{document}
\title{Experiment `p\_13\_10\_recursive\_area\_polygon\_seq' Results}
\author{\today}
\date{}
\maketitle
\textbf{Experiment outcome: }SUCCESS
\\\textbf{Bad responses: }0
\\\textbf{Responses containing}~\texttt{assume}~\textbf{: }0
\\\textbf{Responses containing}~\texttt{expect}~\textbf{: }0
\\\textbf{Resolution attempts: }4
\\\textbf{Hard fails (resolution): }0
\\\textbf{Soft fails (resolution): }0
\\\textbf{Verification attempts: }4
\section*{Problem Specification}
\textbf{Problem name: }p\_13\_10\_recursive\_area\_polygon\_seq
\\\textbf{Natural language statement: }null
\\\textbf{Method signature: }p\_13\_10\_recursive\_area\_polygon\_seq(points: seq<(real, real)>) returns (area: real)
\subsection*{Ensures}
\begin{itemize}
\item \texttt{area == polygon\_area(points)}
\end{itemize}
\subsection*{Requires}
\begin{itemize}
\item \texttt{|points| >= 3}
\end{itemize}
\subsection*{Functional Code Given}
\begin{lstlisting}
function polygon_area(points: seq<(real, real)>): real
  decreases |points|
  requires |points| >= 3
  ensures polygon_area(points) >= 0.0
{
  if |points| == 3 then 
    var (x1, y1) := points[0];
    var (x2, y2) := points[1];
    var (x3, y3) := points[2];
    var sign : real := x1 * (y2 - y3) + x2 * (y3 - y1) + x3 * (y1 - y2);
    if sign < 0.0 then -sign / 2.0 else sign / 2.0 as real 
  else 
    var (x1, y1) := points[0];
    var (x2, y2) := points[1];
    var (x3, y3) := points[2];
    var sign : real := x1 * (y2 - y3) + x2 * (y3 - y1) + x3 * (y1 - y2);
    var triangle_area := if sign < 0.0 then -sign / 2.0 else sign / 2.0 as real;
    triangle_area + polygon_area([points[0]] + points[2..])
}
\end{lstlisting}
\clearpage
\section*{GenAI interactions}
Below you will find all interactions between the `user' (program) and the `assistant' (OpenAI).
\subsection*{Program $\rightarrow$ GenAI}
\begin{lstlisting}
You are given the following task to perform in Dafny:

Write a method to satisfy the below requirements.

The signature should be:

p_13_10_recursive_area_polygon_seq(points: seq<(real, real)>) returns (area: real)

The method should respect the following contract:

requires |points| >= 3, ensures area == polygon_area(points)

The contract uses the following dafny code:

function polygon_area(points: seq<(real, real)>): real
  decreases |points|
  requires |points| >= 3
  ensures polygon_area(points) >= 0.0
{
  if |points| == 3 then 
    var (x1, y1) := points[0];
    var (x2, y2) := points[1];
    var (x3, y3) := points[2];
    var sign : real := x1 * (y2 - y3) + x2 * (y3 - y1) + x3 * (y1 - y2);
    if sign < 0.0 then -sign / 2.0 else sign / 2.0 as real 
  else 
    var (x1, y1) := points[0];
    var (x2, y2) := points[1];
    var (x3, y3) := points[2];
    var sign : real := x1 * (y2 - y3) + x2 * (y3 - y1) + x3 * (y1 - y2);
    var triangle_area := if sign < 0.0 then -sign / 2.0 else sign / 2.0 as real;
    triangle_area + polygon_area([points[0]] + points[2..])
}

These function(s) must not be used in your implementation of the method.

Produce and show only the Dafny body of this method, including the curly braces that surround it. Do not show the signature nor contract. You must not use 'assume' nor 'expect' anywhere in your code.
\end{lstlisting}
\subsection*{GenAI $\rightarrow$ Program}
\textbf{System fingerprint: }null
\\\textbf{ID: }chatcmpl-DnPHTPShQ4RYY6aUWsm0rG57Xlv5C
\\\textbf{Model: }gpt-5.5-2026-04-23
\\\textbf{Created at: }1780667623
\\\textbf{Finish reason: }stop
\\\textbf{Usage: }896 tokens in, and 1367 tokens out
\begin{lstlisting}
{
  if |points| == 3 {
    var (x1, y1) := points[0];
    var (x2, y2) := points[1];
    var (x3, y3) := points[2];
    var sign: real := x1 * (y2 - y3) + x2 * (y3 - y1) + x3 * (y1 - y2);
    area := if sign < 0.0 then -sign / 2.0 else sign / 2.0;
  } else {
    var (x1, y1) := points[0];
    var (x2, y2) := points[1];
    var (x3, y3) := points[2];
    var sign: real := x1 * (y2 - y3) + x2 * (y3 - y1) + x3 * (y1 - y2);
    var triangle_area := if sign < 0.0 then -sign / 2.0 else sign / 2.0;

    var remaining := [points[0]] + points[2..];
    assert |remaining| == |points| - 1;
    assert |remaining| >= 3;

    var rest := p_13_10_recursive_area_polygon_seq(remaining);
    area := triangle_area + rest;
  }
}
\end{lstlisting}
\subsection*{Program $\rightarrow$ GenAI}
\begin{lstlisting}
When using dafny verify, the below error is emitted and verify fails:

problem_data/p_13_10_recursive_area_polygon_seq/program.dfy(24,0): Error: a postcondition could not be proved on this return path
   |
24 | {
   | ^

problem_data/p_13_10_recursive_area_polygon_seq/program.dfy(23,14): Related location: this is the postcondition that could not be proved
   |
23 | 	ensures area == polygon_area(points)
   |               ^^

problem_data/p_13_10_recursive_area_polygon_seq/program.dfy(42,50): Error: cannot prove termination; try supplying a decreases clause
   |
42 |     var rest := p_13_10_recursive_area_polygon_seq(remaining);
   |                                                   ^


Dafny program verifier finished with 2 verified, 2 errors

Correct the error in code by altering only the method body. Produce and show only the Dafny body, including the curly braces that surround it. Do not show the signature nor contract. You must not use 'assume' nor 'expect' anywhere in your code.
\end{lstlisting}
\subsection*{GenAI $\rightarrow$ Program}
\textbf{System fingerprint: }null
\\\textbf{ID: }chatcmpl-DnPILT6Cf7mNf4JCk1Geqb0j6Jgs7
\\\textbf{Model: }gpt-5.5-2026-04-23
\\\textbf{Created at: }1780667677
\\\textbf{Finish reason: }stop
\\\textbf{Usage: }1270 tokens in, and 3541 tokens out
\begin{lstlisting}
{
  var acc: real := 0.0;
  var q := points;

  while |q| > 3
    invariant |q| >= 3
    invariant acc + polygon_area(q) == polygon_area(points)
    decreases |q|
  {
    var (x1, y1) := q[0];
    var (x2, y2) := q[1];
    var (x3, y3) := q[2];
    var sign: real := x1 * (y2 - y3) + x2 * (y3 - y1) + x3 * (y1 - y2);
    var triangle_area: real := if sign < 0.0 then -sign / 2.0 else sign / 2.0;

    var qNext := [q[0]] + q[2..];
    assert |qNext| == |q| - 1;
    assert |qNext| >= 3;
    assert polygon_area(q) == triangle_area + polygon_area(qNext);

    acc := acc + triangle_area;
    q := qNext;
  }

  assert |q| == 3;
  var (x1, y1) := q[0];
  var (x2, y2) := q[1];
  var (x3, y3) := q[2];
  var sign: real := x1 * (y2 - y3) + x2 * (y3 - y1) + x3 * (y1 - y2);
  var triangle_area: real := if sign < 0.0 then -sign / 2.0 else sign / 2.0;
  assert polygon_area(q) == triangle_area;

  area := acc + triangle_area;
  assert area == polygon_area(points);
}
\end{lstlisting}
\subsection*{Program $\rightarrow$ GenAI}
\begin{lstlisting}
When using dafny verify, the below error is emitted and verify fails:

problem_data/p_13_10_recursive_area_polygon_seq/program.dfy(42,27): Error: assertion might not hold
   |
42 |     assert polygon_area(q) == triangle_area + polygon_area(qNext);
   |                            ^^


Dafny program verifier finished with 2 verified, 1 error

Correct the error in code by altering only the method body. Produce and show only the Dafny body, including the curly braces that surround it. Do not show the signature nor contract. You must not use 'assume' nor 'expect' anywhere in your code.
\end{lstlisting}
\subsection*{GenAI $\rightarrow$ Program}
\textbf{System fingerprint: }null
\\\textbf{ID: }chatcmpl-DnPK5KTNdS5kUiYwmGFwhz00VAAFE
\\\textbf{Model: }gpt-5.5-2026-04-23
\\\textbf{Created at: }1780667785
\\\textbf{Finish reason: }stop
\\\textbf{Usage: }1745 tokens in, and 4051 tokens out
\begin{lstlisting}
{
  var acc: real := 0.0;
  var q := points;
  assert q == points;
  assert acc + polygon_area(q) == polygon_area(points);

  while |q| > 3
    invariant |q| >= 3
    invariant acc + polygon_area(q) == polygon_area(points)
    decreases |q|
  {
    var (x1, y1) := q[0];
    var (x2, y2) := q[1];
    var (x3, y3) := q[2];
    var sign: real := x1 * (y2 - y3) + x2 * (y3 - y1) + x3 * (y1 - y2);
    var triangle_area: real := if sign < 0.0 then -sign / 2.0 else sign / 2.0;

    var qNext := [q[0]] + q[2..];
    assert 4 <= |q|;
    assert |q[2..]| == |q| - 2;
    assert |qNext| == |q| - 1;
    assert |qNext| >= 3;
    assert qNext == [q[0]] + q[2..];
    assert q[0] == (x1, y1);
    assert q[1] == (x2, y2);
    assert q[2] == (x3, y3);

    calc {
      polygon_area(q);
    == {
      assert |q| != 3;
      assert q[0] == (x1, y1);
      assert q[1] == (x2, y2);
      assert q[2] == (x3, y3);
    }
      triangle_area + polygon_area([q[0]] + q[2..]);
    == {
      assert qNext == [q[0]] + q[2..];
    }
      triangle_area + polygon_area(qNext);
    }

    calc {
      (acc + triangle_area) + polygon_area(qNext);
    == {
    }
      acc + (triangle_area + polygon_area(qNext));
    == {
      assert polygon_area(q) == triangle_area + polygon_area(qNext);
    }
      acc + polygon_area(q);
    == {
      assert acc + polygon_area(q) == polygon_area(points);
    }
      polygon_area(points);
    }

    acc := acc + triangle_area;
    q := qNext;
  }

  assert |q| == 3;
  var (x1, y1) := q[0];
  var (x2, y2) := q[1];
  var (x3, y3) := q[2];
  var sign: real := x1 * (y2 - y3) + x2 * (y3 - y1) + x3 * (y1 - y2);
  var triangle_area: real := if sign < 0.0 then -sign / 2.0 else sign / 2.0;
  assert q[0] == (x1, y1);
  assert q[1] == (x2, y2);
  assert q[2] == (x3, y3);

  calc {
    polygon_area(q);
  == {
    assert |q| == 3;
    assert q[0] == (x1, y1);
    assert q[1] == (x2, y2);
    assert q[2] == (x3, y3);
  }
    triangle_area;
  }

  area := acc + triangle_area;
  calc {
    area;
  == {
  }
    acc + triangle_area;
  == {
    assert polygon_area(q) == triangle_area;
  }
    acc + polygon_area(q);
  == {
    assert acc + polygon_area(q) == polygon_area(points);
  }
    polygon_area(points);
  }
}
\end{lstlisting}
\subsection*{Program $\rightarrow$ GenAI}
\begin{lstlisting}
When using dafny verify, the below error is emitted and verify fails:

problem_data/p_13_10_recursive_area_polygon_seq/program.dfy(59,20): Error: the calculation step between the previous line and this line could not be proved
   |
59 |       triangle_area + polygon_area([q[0]] + q[2..]);
   |                     ^


Dafny program verifier finished with 2 verified, 1 error

Correct the error in code by altering only the method body. Produce and show only the Dafny body, including the curly braces that surround it. Do not show the signature nor contract. You must not use 'assume' nor 'expect' anywhere in your code.
\end{lstlisting}
\subsection*{GenAI $\rightarrow$ Program}
\textbf{System fingerprint: }null
\\\textbf{ID: }chatcmpl-DnPLnw7HDd7LUiurYOJapoCNm187L
\\\textbf{Model: }gpt-5.5-2026-04-23
\\\textbf{Created at: }1780667891
\\\textbf{Finish reason: }stop
\\\textbf{Usage: }2853 tokens in, and 7905 tokens out
\begin{lstlisting}
{
  var acc: real := 0.0;
  var q := points;
  assert acc + polygon_area(q) == polygon_area(points);

  while |q| > 3
    invariant |q| >= 3
    invariant acc + polygon_area(q) == polygon_area(points)
    decreases |q|
  {
    var sign: real := (q[0]).0 * ((q[1]).1 - (q[2]).1) + (q[1]).0 * ((q[2]).1 - (q[0]).1) + (q[2]).0 * ((q[0]).1 - (q[1]).1);
    var triangle_area: real := if sign < 0.0 then -sign / 2.0 else sign / 2.0;

    var qNext := [q[0]] + q[2..];
    assert |q[2..]| == |q| - 2;
    assert |qNext| == |q| - 1;
    assert |qNext| >= 3;
    assert qNext == [q[0]] + q[2..];

    assert polygon_area(q) == triangle_area + polygon_area(qNext) by {
      reveal polygon_area();
      assert |q| != 3;
      assert qNext == [q[0]] + q[2..];
    }

    var newAcc := acc + triangle_area;
    assert newAcc + polygon_area(qNext) == acc + (triangle_area + polygon_area(qNext));
    assert newAcc + polygon_area(qNext) == acc + polygon_area(q);
    assert newAcc + polygon_area(qNext) == polygon_area(points);

    acc := newAcc;
    q := qNext;
  }

  assert |q| == 3;
  var sign: real := (q[0]).0 * ((q[1]).1 - (q[2]).1) + (q[1]).0 * ((q[2]).1 - (q[0]).1) + (q[2]).0 * ((q[0]).1 - (q[1]).1);
  var triangle_area: real := if sign < 0.0 then -sign / 2.0 else sign / 2.0;

  assert polygon_area(q) == triangle_area by {
    reveal polygon_area();
    assert |q| == 3;
  }

  area := acc + triangle_area;
  assert area == acc + polygon_area(q);
  assert area == polygon_area(points);
}
\end{lstlisting}
\section*{Final Program}
\begin{lstlisting}
function polygon_area(points: seq<(real, real)>): real
  decreases |points|
  requires |points| >= 3
  ensures polygon_area(points) >= 0.0
{
  if |points| == 3 then 
    var (x1, y1) := points[0];
    var (x2, y2) := points[1];
    var (x3, y3) := points[2];
    var sign : real := x1 * (y2 - y3) + x2 * (y3 - y1) + x3 * (y1 - y2);
    if sign < 0.0 then -sign / 2.0 else sign / 2.0 as real 
  else 
    var (x1, y1) := points[0];
    var (x2, y2) := points[1];
    var (x3, y3) := points[2];
    var sign : real := x1 * (y2 - y3) + x2 * (y3 - y1) + x3 * (y1 - y2);
    var triangle_area := if sign < 0.0 then -sign / 2.0 else sign / 2.0 as real;
    triangle_area + polygon_area([points[0]] + points[2..])
}

method p_13_10_recursive_area_polygon_seq(points: seq<(real, real)>) returns (area: real)
	requires |points| >= 3
	ensures area == polygon_area(points)
{
  var acc: real := 0.0;
  var q := points;
  assert acc + polygon_area(q) == polygon_area(points);

  while |q| > 3
    invariant |q| >= 3
    invariant acc + polygon_area(q) == polygon_area(points)
    decreases |q|
  {
    var sign: real := (q[0]).0 * ((q[1]).1 - (q[2]).1) + (q[1]).0 * ((q[2]).1 - (q[0]).1) + (q[2]).0 * ((q[0]).1 - (q[1]).1);
    var triangle_area: real := if sign < 0.0 then -sign / 2.0 else sign / 2.0;

    var qNext := [q[0]] + q[2..];
    assert |q[2..]| == |q| - 2;
    assert |qNext| == |q| - 1;
    assert |qNext| >= 3;
    assert qNext == [q[0]] + q[2..];

    assert polygon_area(q) == triangle_area + polygon_area(qNext) by {
      reveal polygon_area();
      assert |q| != 3;
      assert qNext == [q[0]] + q[2..];
    }

    var newAcc := acc + triangle_area;
    assert newAcc + polygon_area(qNext) == acc + (triangle_area + polygon_area(qNext));
    assert newAcc + polygon_area(qNext) == acc + polygon_area(q);
    assert newAcc + polygon_area(qNext) == polygon_area(points);

    acc := newAcc;
    q := qNext;
  }

  assert |q| == 3;
  var sign: real := (q[0]).0 * ((q[1]).1 - (q[2]).1) + (q[1]).0 * ((q[2]).1 - (q[0]).1) + (q[2]).0 * ((q[0]).1 - (q[1]).1);
  var triangle_area: real := if sign < 0.0 then -sign / 2.0 else sign / 2.0;

  assert polygon_area(q) == triangle_area by {
    reveal polygon_area();
    assert |q| == 3;
  }

  area := acc + triangle_area;
  assert area == acc + polygon_area(q);
  assert area == polygon_area(points);
}
\end{lstlisting}
\section*{Total Token Usage}
\textbf{Input tokens: }6764
\\\textbf{Output tokens: }16864
\\\textbf{Reasoning tokens: }14386
\\\textbf{Sum of `total tokens': }23628
\section*{Experiment Timings}
\textbf{Overall Experiment} started at 1780667623432, ended at 1780668070384, lasting 446952ms (446.95 seconds)
\\\textbf{Iteration \#4} started at 1780667890319, ended at 1780668070384, lasting 180065ms (180.07 seconds)
\\\textbf{Iteration \#1} started at 1780667623433, ended at 1780667676577, lasting 53144ms (53.14 seconds)
\\\textbf{Iteration \#2} started at 1780667676580, ended at 1780667784433, lasting 107853ms (107.85 seconds)
\\\textbf{Iteration \#3} started at 1780667784433, ended at 1780667890298, lasting 105865ms (105.87 seconds)
\end{document}
